% GENERATED FILE. Edit canonical lessons, then run npm run generate:books.
% canonical-source-hash: fnv1a64:2903cfdc40655620
% canonical-lessons: AR-C362-mantiq, AR-C362-mughalata, AR-C362-aql, AR-C362-istantaja, AR-C362-tanaqud

\chapter{Logic and the mind}
\label{ch:ar-logic-and-the-mind}
\hlchaptermodality{362}

\begin{chapteropening}
\textbf{By the end of this chapter:} \emph{I can talk about logic, a fallacy, the mind and a contradiction, and say what someone concluded.}

Complete the last lesson of chapter 362: I can talk about logic, a fallacy, the mind and a contradiction, and say what someone concluded.
\end{chapteropening}

\section[\texorpdfstring{\ar{منطق} (manṭiq)}{منطق (manṭiq)}]{\ar{منطق} (manṭiq) — logic}

\label{lesson:AR-C362-mantiq}



\begin{quote}
Before the new one: say the Arabic for to stir up, then the Arabic for to wonder.
\end{quote}

\subsection*{You'll want to know: \ar{منطق}}
\textbf{\ar{منطق}} — \emph{manṭiq} — \textquotedblleft{}logic\textquotedblright{}.

\begin{quote}
\ar{لا} \ar{أفهم} \ar{منطق} \ar{هذا} \ar{القرار}.
\end{quote}

\begin{quote}
\emph{lā afhamu manṭiq hādhā al-qarār.}
\end{quote}

\begin{quote}
I do not understand the logic of this decision.
\end{quote}

The adjective is \ar{منطقي}, \emph{manṭiqī}, logical.

\begin{cousinweb}
\textbf{\ar{منطق}} (\emph{manṭiq}) is built on the native Arabic root \emph{n-ṭ-q}, \textquotedblleft{}to speak\textquotedblright{} (Lane). Lane gives \emph{manṭiq} as speech, the putting of ideas into words; later it became the word for logic.
\end{cousinweb}

\begin{grammarlens}[title={What you've built}]
One more word for giving reasons and opinions.
\end{grammarlens}

\subsection*{Your turn}
\begin{itemize}
\raggedright
  \item \emph{Say it:} \emph{manṭiq}
  \item \emph{Say it:} \emph{manṭiq}, once more
  \item \emph{Say it:} \emph{lā afhamu manṭiq hādhā al-qarār.} — I do not understand the logic of this decision.
  \item \emph{Recall:} say the Arabic for to stir up, then the Arabic for to wonder, then say \emph{manṭiq} again
\end{itemize}

\begin{tcolorbox}[breakable,colback=teal!4,colframe=teal!35!black,title={Before you move on}]
What does \ar{منطق} mean? (\textquotedblleft{}Logic\textquotedblright{}.) Say it once more.
\end{tcolorbox}
\section[\texorpdfstring{\ar{مغالطة} (mughālaṭa)}{مغالطة (mughālaṭa)}]{\ar{مغالطة} (mughālaṭa) — a fallacy}

\label{lesson:AR-C362-mughalata}



\begin{quote}
Before the new one: say the Arabic for to wonder, then the Arabic for logic.
\end{quote}

\subsection*{You'll want to know: \ar{مغالطة}}
\textbf{\ar{مغالطة}} — \emph{mughālaṭa} — \textquotedblleft{}a fallacy\textquotedblright{}.

\begin{quote}
\ar{في} \ar{هذا} \ar{الكلام} \ar{مغالطة} \ar{واضحة}.
\end{quote}

\begin{quote}
\emph{fī hādhā al-kalām mughālaṭa wāḍiḥa.}
\end{quote}

\begin{quote}
There is a clear fallacy in what was said.
\end{quote}

\begin{cousinweb}
\textbf{\ar{مغالطة}} (\emph{mughālaṭa}) is built on the native Arabic root \emph{gh-l-ṭ}, \textquotedblleft{}to make a mistake\textquotedblright{} (Lane). It is the Form III verbal noun, leading someone into error: a fallacy, an argument that misleads.
\end{cousinweb}

\begin{grammarlens}[title={What you've built}]
One more word for giving reasons and opinions.
\end{grammarlens}

\subsection*{Your turn}
\begin{itemize}
\raggedright
  \item \emph{Say it:} \emph{mughālaṭa}
  \item \emph{Say it:} \emph{mughālaṭa}, once more
  \item \emph{Say it:} \emph{fī hādhā al-kalām mughālaṭa wāḍiḥa.} — There is a clear fallacy in what was said.
  \item \emph{Recall:} say the Arabic for to wonder, then the Arabic for logic, then say \emph{mughālaṭa} again
\end{itemize}

\begin{tcolorbox}[breakable,colback=teal!4,colframe=teal!35!black,title={Before you move on}]
What does \ar{مغالطة} mean? (\textquotedblleft{}A fallacy\textquotedblright{}.) Say it once more.
\end{tcolorbox}
\section[\texorpdfstring{\ar{عقل} (ʿaql)}{عقل (ʿaql)}]{\ar{عقل} (ʿaql) — the mind}

\label{lesson:AR-C362-aql}



\begin{quote}
Before the new one: say the Arabic for logic, then the Arabic for a fallacy.
\end{quote}

\subsection*{You'll want to know: \ar{عقل}}
\textbf{\ar{عقل}} — \emph{ʿaql} — \textquotedblleft{}the mind\textquotedblright{}.

\begin{quote}
\ar{فكر} \ar{الرجل} \ar{بعقله} \ar{لا} \ar{بقلبه}.
\end{quote}

\begin{quote}
\emph{fakkara ar-rajul bi-ʿaqlihi lā bi-qalbih.}
\end{quote}

\begin{quote}
The man thought with his head, not with his heart.
\end{quote}

Plural: \ar{عقول}, \emph{ʿuqūl}.

\begin{cousinweb}
\textbf{\ar{عقل}} (\emph{ʿaql}) is built on the native Arabic root \emph{ʿ-q-l}, \textquotedblleft{}to tie, to restrain\textquotedblright{} (Lane). Lane gives the root as tying a camel's leg; \emph{ʿaql} is the mind, which holds a person back from folly.
\end{cousinweb}

\begin{grammarlens}[title={What you've built}]
One more word for giving reasons and opinions.
\end{grammarlens}

\subsection*{Your turn}
\begin{itemize}
\raggedright
  \item \emph{Say it:} \emph{ʿaql}
  \item \emph{Say it:} \emph{ʿaql}, once more
  \item \emph{Say it:} \emph{fakkara ar-rajul bi-ʿaqlihi lā bi-qalbih.} — The man thought with his head, not with his heart.
  \item \emph{Recall:} say the Arabic for logic, then the Arabic for a fallacy, then say \emph{ʿaql} again
\end{itemize}

\begin{tcolorbox}[breakable,colback=teal!4,colframe=teal!35!black,title={Before you move on}]
What does \ar{عقل} mean? (\textquotedblleft{}The mind\textquotedblright{}.) Say it once more.
\end{tcolorbox}
\section[\texorpdfstring{\ar{استنتج} (istantaja)}{استنتج (istantaja)}]{\ar{استنتج} (istantaja) — to deduce; to conclude}

\label{lesson:AR-C362-istantaja}



\begin{quote}
Before the new one: say the Arabic for a fallacy, then the Arabic for the mind.
\end{quote}

\subsection*{You'll want to know: \ar{استنتج}}
\textbf{\ar{استنتج}} — \emph{istantaja} — \textquotedblleft{}to deduce; to conclude\textquotedblright{}.

\begin{quote}
\ar{استنتجت} \ar{من} \ar{كلامه} \ar{أنه} \ar{متعب}.
\end{quote}

\begin{quote}
\emph{istantajtu min kalāmihi annahu mutʿab.}
\end{quote}

\begin{quote}
I concluded from what he said that he was tired.
\end{quote}

Present: \ar{يستنتج}, \emph{yastantij}.

Followed by \ar{أن}: \ar{استنتج} \ar{أن} …. A related noun is \ar{نتيجة}, a result.

\begin{cousinweb}
\textbf{\ar{استنتج}} (\emph{istantaja}) is built on the native Arabic root \emph{n-t-j}, \textquotedblleft{}to bring forth young\textquotedblright{} (Lane). It is Form X, to seek the outcome: to draw a conclusion.
\end{cousinweb}

\begin{grammarlens}[title={What you've built}]
One more word for giving reasons and opinions.
\end{grammarlens}

\subsection*{Your turn}
\begin{itemize}
\raggedright
  \item \emph{Say it:} \emph{istantaja}
  \item \emph{Say it:} \emph{istantaja}, once more
  \item \emph{Say it:} \emph{istantajtu min kalāmihi annahu mutʿab.} — I concluded from what he said that he was tired.
  \item \emph{Recall:} say the Arabic for a fallacy, then the Arabic for the mind, then say \emph{istantaja} again
\end{itemize}

\begin{tcolorbox}[breakable,colback=teal!4,colframe=teal!35!black,title={Before you move on}]
What does \ar{استنتج} mean? (\textquotedblleft{}To deduce; to conclude\textquotedblright{}.) Say it once more.
\end{tcolorbox}
\section[\texorpdfstring{\ar{تناقض} (tanāquḍ)}{تناقض (tanāquḍ)}]{\ar{تناقض} (tanāquḍ) — a contradiction}

\label{lesson:AR-C362-tanaqud}



\begin{quote}
Before the new one: say the Arabic for the mind, then the Arabic for to deduce.
\end{quote}

\subsection*{You'll want to know: \ar{تناقض}}
\textbf{\ar{تناقض}} — \emph{tanāquḍ} — \textquotedblleft{}a contradiction\textquotedblright{}.

\begin{quote}
\ar{هناك} \ar{تناقض} \ar{بين} \ar{كلامه} \ar{وعمله}.
\end{quote}

\begin{quote}
\emph{hunāka tanāquḍ bayna kalāmihi wa-ʿamalih.}
\end{quote}

\begin{quote}
There is a contradiction between his words and his deeds.
\end{quote}

\begin{cousinweb}
\textbf{\ar{تناقض}} (\emph{tanāquḍ}) is built on the native Arabic root \emph{n-q-ḍ}, \textquotedblleft{}to undo, to unravel\textquotedblright{} (Lane). It is the Form VI verbal noun: two things that undo each other, a contradiction.
\end{cousinweb}

\begin{grammarlens}[title={What you've built}]
One more word for giving reasons and opinions.
\end{grammarlens}

\subsection*{Your turn}
\begin{itemize}
\raggedright
  \item \emph{Say it:} \emph{tanāquḍ}
  \item \emph{Say it:} \emph{tanāquḍ}, once more
  \item \emph{Say it:} \emph{hunāka tanāquḍ bayna kalāmihi wa-ʿamalih.} — There is a contradiction between his words and his deeds.
  \item \emph{Recall:} say the Arabic for the mind, then the Arabic for to deduce, then say \emph{tanāquḍ} again
\end{itemize}

\begin{tcolorbox}[breakable,colback=teal!4,colframe=teal!35!black,title={Before you move on}]
What does \ar{تناقض} mean? (\textquotedblleft{}A contradiction\textquotedblright{}.) Say it once more.
\end{tcolorbox}
